% Abhakseite „Das kann ich“ – Zone nur im Gesamt (\mitzone)
%% TODO Ich-kann-Sätze: Platzhalter wie in den Aufgabendateien
\begin{abhakseite}
\ifdefined\mitzone
\abhakgruppe{Kennst du schon}
\abhak{1}{Verbindungsvektor bilden (Spitze minus Fuß) und den Betrag als Wurzel aus der Quadratsumme berechnen}
\abhak{2}{Wurzel- und Betragsgleichungen lösen (quadrieren, beide Vorzeichen ansetzen, Lösungen am Sachzusammenhang sieben)}
\abhak{3}{Koordinatengleichung einer Ebene lesen und den Normalenvektor ablesen}
\abhak{4}{Geradengleichung in Parameterform lesen und einen Laufpunkt mit Parameter ansetzen}
\abhak{5}{Skalarprodukt berechnen und „null heißt senkrecht“ anwenden}
\abhak{6}{Satz des Pythagoras, Satz des Thales und Strahlensätze als Begründungsfiguren}
\abhak{7}{Verbindungsvektor bilden (Spitze minus Fuß) und den Betrag als Wurzel aus der Quadratsumme berechnen – Fehler finden}
\abhak{8}{Verbindungsvektor bilden (Spitze minus Fuß) und den Betrag als Wurzel aus der Quadratsumme berechnen – selbst rechnen}
\fi
\abhakgruppe{Einheit 1 · Abstand zweier Punkte}
\abhak{9}{Hin oder zurück?}
\abhak{10}{Abstand zweier Punkte}
\abhak{11}{Abstand zweier Punkte – Fehler finden, Begründen, Anwendung}
\abhakgruppe{Einheit 2 · Abstand Punkt–Ebene}
\abhak{12}{Abstand Punkt–Ebene}
\abhak{13}{Abstand Punkt–Ebene – Fehler finden, Begründen, Anwendung}
\abhakgruppe{Einheit 3 · Lotfußpunkt}
\abhak{14}{Lotfußpunkt}
\abhak{15}{Lotfußpunkt – Fehler finden, Begründen, Anwendung}
\abhakgruppe{Einheit 4 · Abstand als Argument}
\abhak{16}{Abstand als Argument}
\abhak{17}{Abstand als Argument – Fehler finden, Begründen, Anwendung}
\end{abhakseite}
